Figure~\ref{fig:atlas-conditional-events} visualizes conditioning as restricting attention to a specified event.

\begin{figure}[!htbp]
\centering
\begin{tikzpicture}[scale=1.45,>=Stealth]
\draw[black!60] (-2.5,-1.6) rectangle (2.5,1.6);
\fill[figblue!15] (0.65,0) circle (1.25);
\begin{scope}
\clip (-0.65,0) circle (1.25);
\fill[figgreen!45] (0.65,0) circle (1.25);
\end{scope}
\draw[figred,thick] (-0.65,0) circle (1.25);
\draw[figblue,thick] (0.65,0) circle (1.25);
\node at (-1.25,0.25) {$A$};
\node at (1.25,0.25) {$B$};
\node at (0,0) {$A\cap B$};
\node at (2.2,1.3) {$\Omega$};
\end{tikzpicture}
\caption{The outer rectangle represents the sample space. Conditioning on $B$ restricts attention to the blue disk; its overlap with $A$ is highlighted in green. For $\Prob(B)>0$, the conditional probability is $\Prob(A\mid B)=\Prob(A\cap B)/\Prob(B)$. Areas are schematic and need not be proportional to probabilities.}
\label{fig:atlas-conditional-events}
\end{figure}